\documentclass{article}
\usepackage[T1]{fontenc}
\usepackage{lmodern}
\usepackage{graphicx}
\usepackage{amsmath,amssymb}
\usepackage[a4paper,margin=2cm]{geometry}
\usepackage[absolute,overlay]{textpos}
\setlength{\TPHorizModule}{1mm}
\setlength{\TPVertModule}{1mm}
\usepackage[indonesian]{babel}
\usepackage{hyperref}
\setlength{\emergencystretch}{2em}
% unit: Halaman 19

\begin{document}

% === Halaman 19 (python/native) ===

(c) Dekripsi  (oleh Alice)

Alice mendekripsi cipherteks dari Bob dengan kunci privatnya (x = 243, p = 2273):

    (i) Dekripsi c1 = (1439, 74)

(ax)– 1= ap – 1 – x mod p = 14392273 – 1 - 243 mod 2273 = 14392029 mod 2273 = 1791

m1 = b/ax mod p = b(ax)–1mod p = 74 1791 mod 2273 = 700 = 0700

    (ii) Dekripsi c2 = (1220, 1682)

(ax)– 1= ap – 1 – x mod p = 12202029 mod 2273 = 1125

m2 = b/ax mod p = 1682(ax)–1mod p = 1682 1125 mod 2273 = 1114

Plainteks yang didekripsi adalah m1m2 = 07001114, yang kalau dikodekan menjadi teks 

adalah empat digit adalah “HALO”, sama dengan plainteks yang dikirim oleh Bob. 

19

\end{document}
